% 从旧数学附录迁出的连通性与动态树说明；不新增算法。
\section{SCC 模板题与速度榜口径}\label{knowledge-scc-template}
\href{https://qoj.ac/problem/906}{QOJ906} 来自 Library Checker，要求输出强连通分量的拓扑序。
本库 Tarjan 编号按逆拓扑序产生，驱动将输入点号加一，再按分量编号从大到小输出，输出点号减一。
\href{https://qoj.ac/submission/2940638}{CCF\_NOI 的记录 2940638} 通过全部 11 个测试，203 ms、58636 KB。
验证覆盖 add、run、bel 和完整分量输出，不调用 dag。
本地独立可达性参考核对分量等价关系及每条跨分量边方向，并测试 50 万点链、环和多分量。
普通与 ASan/UBSan 均通过；本地测试线程使用 1 GiB 栈容纳 sanitizer 扩大的递归栈帧，
提交代码保持普通 main 和递归 DFS。

2026-09-12 核对速度榜时，页面第 66 行处于 203 ms 的第 65--67 行并列区间。
但该表合并同一用户的提交，不能将行号与满分提交数 578 配对。
以 CCF\_NOI 逐页读取页面链接的满分列表，共取得 143 条，其中 93 条严格更快、3 条同为 203 ms。
可见范围内按耗时并列排名为 94 out of 143，位置区间为 94--96。
统计页显示所有提交 1053、满分提交 578；与可见列表的计数差异原因尚未核实，因此全站排名仍为待查。
记录清单与证据见 \texttt{docs/ranking-audits.json}，不得将该局部排名改写为全站排名。

\section{并查集模板题的接口范围}\label{knowledge-dsu-template}
\href{https://www.luogu.com.cn/problem/P3367}{P3367} 输入点号为 1-based，驱动统一减一后调用 dsu。
\href{https://www.luogu.com.cn/record/297668102}{记录 297668102} 通过全部 20 个测试，最慢单点 196 ms、峰值内存 2.27 MB。
该输出验证 merge、same 及其调用的 find，不验证 size、groups；后两者已由独立标签参考和 50 万点长链测试核对。

\section{ACL 的组合约定}\label{knowledge-acl-compose}
固定参考 AtCoder Library 提交 \texttt{864245a00b00}，公开模块分为数据结构、图论、数学与字符串。
学习短小的运算接口及明确的组合契约，不直接复制全部工程支持层。
尤其注意：懒标记 composition(f,g) 表示先作用 g、再作用 f；SCC 分量方向与 2-SAT 的真假编码不能分别替换。
本库已补齐 ModInt 的四种复合赋值，Fenwick 的模数区间和组合通过独立整数余数参考及 ASan/UBSan。新增运算只有本地验证，不能套用历史 OJ AC；取模和仍不适用于 kth。
完整 12 模块、6 内部头文件的对照及缺项清单见 \texttt{docs/ACL-REVIEW.md}。

\section{动态树路径乘积：虚子树与双向信息}\label{knowledge-lct-product}\index{动态树 DP}\index{LCT!虚子树}
附件“爱莲说”对应的量是
\[
F(u)=\sum_{v\in V}\prod_{x\in\operatorname{path}(u,v)}w_x.
\]
包括单点路径，所以一个孤立点的答案是其点权，而不是 1。
固定根后，递推为 $F(u)=w_u(1+\sum_{v\text{ 为子节点}}F(v))$。

可以单点修改并动态选取查询根。它维护固定树，不包含连边、删边。

辅助树按实链顺序维护三项：整条链点权积 $P$、从链首出发的贡献 $L$、
从链尾出发的贡献 $R$。每个点另记虚儿子贡献之和 $G$。
对于左、右辅助子树，合并式为
\[
P=P_lwP_r,\qquad
L=L_l+P_lw(1+G+L_r),\qquad
R=R_r+P_rw(1+G+R_l).
\]
空辅助树取 $P=1,L=R=0$。
翻转时必须同时交换左右儿子和 $L,R$；只翻儿子会破坏换根后的方向。

\texttt{access} 把旧右链变为虚儿子，将它的 $L$ 加入 $G$；
把新右链变为实儿子，从 $G$ 减去它的 $L$。旧右链的父指针保留为路径父亲。
修改点权前先将该点换为根，确保旧贡献已经沿访问路径正确移出虚子树项；
然后修改权值并更新合并信息。换根结束时根处的 $L$ 就是所求总和。
整个过程只用加减乘，没有除法，零权、负权与合数模数同样适用。

接口使用 0-based 顶点；内部辅助树为 1-based，内部函数不可直接作为外部 API。
附件输入参数 $x$ 对应 $w=1-x$，这个转换在使用程序中显式完成。
从源码可确定上述递推和输入协议，但原题外部地址尚未核验，因此只登记本地应用用法。
建树 $O(n)$，修改和查询摊还 $O(\log n)$，总空间 $O(n)$；
建树和辅助树下传保持递归，最坏深度 $O(n)$。
本地以逐终点路径枚举、模 8、零权、20 万点链与星形及反复换根验证，
不将普通 LCT 的线上记录视为这个变体的 AC。

\section{隐式 Treap 位翻转与顺序反转}\label{knowledge-treap-flip}\index{Treap!位翻转}
附件 FHQ Treap 中名为 \texttt{Reverse} 的操作是逐位异或 1；
它没有交换左右儿子，不是区间顺序反转。

\texttt{flip\_bits(l,r)} 明确表达前者，要求区间元素均为 0 或 1；
原有 \texttt{reverse(l,r)} 表达后者。两者均使用 1-based 闭区间。
例如 0011 的前三位取反得到 1101，而前三位顺序反转得到 1001。
二进制序列的区间和就是 1 的个数，仍可直接调用 \texttt{query}。

取反可写作 $x\mapsto1-x$。用符号标记与加法标记表示 $x\mapsto sx+b$，
其中 $s\in\{1,-1\}$。取负时同时令 $s\gets-s,b\gets-b$，再加 1；
下传时必须先符号、再加法。若把先后顺序颠倒，已有加法会变成错误的符号。
顺序反转则只交换儿子并翻转其标记，它与逐元素的取负、加法可交换。
结构体公开位翻转接口仅用于二进制值；所有标记和取负中间量仍须在
\texttt{long long} 范围内，不能对 \texttt{LLONG\_MIN} 直接取负。

附件 \texttt{Build} 按中点建立树后给节点随机优先级，未修复优先级堆序；
这不能直接继承随机 Treap 的平衡复杂度证明。
本库逐次 \texttt{insert} 构建，期望 $O(n\log n)$，每次分裂合并均维持堆序；
不能把这种构建成本写为线性。随机优先级下单次操作期望 $O(\log n)$，
不保证针对已知种子的自适应对抗操作的最坏界。分裂、合并及中序输出保持递归。
附件缺少完整牛客题号和主程序，因此新增使用例只作为本地接口示例。

\section{替罪羊树：重复计数与墓碑清理}\label{knowledge-scapegoat}\index{替罪羊树}

\texttt{size} 计元素副本总数，\texttt{nodes} 计物理节点数，
\texttt{live} 计仍有副本的不同键数。重复值增加计数，删除最后一个副本后保留墓碑。
查询排名只累计副本数，不能用物理节点数代替。

更新回溯时，若某孩子占节点数的比例超过 $3/4$，或者有效键不足物理节点的一半，
就递归中序收集活节点，按中点重建。重建复用活节点编号，保留重复次数。
重建后的中序和计数均不变，故集合与排名不变；所有新子树均满足平衡与有效比例。
令 $d$ 为当前不同键数，树高为 $O(\log(d+2))$，因为每层节点数缩小到至多 $3/4$，
根的物理节点数至多为 $2d$。查询最坏对数，更新摊还对数；一次重建仍可能线性。
删除清理后的墓碑不回收到节点池，空间按累计分配量计算，不按当前集合大小计算。

附件的 \texttt{Rank} 只有命中相等键才加一，未命中时可能少一；
仅插入 1 后查询 2 的排名就是一个反例。
本版统一定义为严格较小副本数加一，空树中的排名也为 1。
严格前驱后继使用 \texttt{optional}，避免数值哨兵与合法键冲突；
完整 \texttt{long long} 键无需做减一或加一。
P3369 保证第 k 小及邻居存在，使用程序才能直接解引用结果。
所有递归保持原样；本地排序向量对拍、重建不变量及 20 万键规模均已验证，线上待补。

